\begin{proof}[The specified Hopf gluing and the cited surgery calculation]
\begin{enumerate}[label=\arabic*.]
\item Orient the Hopf circle bundle by the circle followed by the positive
two-sphere, and use the outward-normal-first orientation of
$\partial(D^4_u\times D^4_y)$. Our gluing is
\[
 \lambda(f)=(D^4_u\times S^3_y)\cup_{\widetilde f}(S^3_u\times D^4_y),
 \qquad x\sim\widetilde f(x),
\]
where $\widetilde f$ is a principal-circle lift. Put
$K=S^2\times D^4$ and $N_f=K_+\cup_f(-K_-)$. This is the marked double
in \cite[Lemma~6.2, pp.~34--35]{Jin2026}. Its integral cohomology is
$\mathbb Z$ in degrees $0,2,4,6$ and zero otherwise; it is simply
connected and spin. Let $x$ restrict to the positive generator on both
cores, and choose $y$ by $\langle xy,[N_f]\rangle=1$. The cited collar
bordism is obtained from $K\times[0,1]_t$ by identifying the two inward
collars through $(z,s,0)\sim(f(z),s,1)$. Its boundary is the reverse-seam
mapping torus together with $-N_f$: contraction of
$\operatorname{or}(S^2\times S^3)\,ds\,dt$ by $-\partial_s$ gives
$\operatorname{or}(S^2\times S^3)\,dt$ on the mapping-torus boundary.
The marked class and preferred spin structure extend. Thus
\[
 x^2=k(f)y,\qquad p_1(TN_f)=\ell(f)y. \tag{P3F.22.2}
\]

\item Set $f_0=f\circ a^{-k(f)}$. The explicit quaternionic extension
of \cref{night:P3F.4}, on $D^4_u\times S^3_y$, together with the identity
on the other gluing piece, gives an orientation-preserving
$\lambda(f_0)\cong\lambda(f)$. Jin's additive coordinates give
$k(f_0)=0$ and $p:=\ell(f_0)=\ell(f)-4k(f)=-24A(f)$.
This uses the actual extension and does not assume additivity of
$\lambda$.

\item Let $L\to N_{f_0}$ have $c_1(L)=x$. The circle bundle decomposes as
\[
 S(L)=(S^3\times D^4)_+\cup_{\widetilde f_0}-(S^3\times D^4)_-.
\]
Attach a four-handle $D^4\times D^4$ to $D(L)$ along the first
$S^3\times D^4$, using its product framing. Its boundary replaces this
piece by $-(D^4\times S^3)$: the boundary orientations of
$S^3\times D^4$ and $D^4\times S^3$ are respectively the negative and
positive product orientations. Consequently the surgery boundary is
exactly $-\lambda(f_0)$.

This is Jiang's disc-bundle and four-handle construction, with the exact
substitution
\[
 [N,t]=(N_{f_0},x),\qquad
 (k,p,\varepsilon,\delta)=(0,\ell(f)-4k(f),0,0),
\]
where $\varepsilon$ and $\delta$ denote the source's second
Stiefel--Whitney and Kirby--Siebenmann indicators. The hypotheses hold
because $N_{f_0}$ is smooth and spin and $x^2=0$. We use the calculated
intersection and characteristic numbers from
\cite[Lemma~3.5, Case~2, pp.~13--15]{Jiang2014}.
In its alternative description the Hopf-disc neighbourhood is replaced
by $\mathbb CP^2\times D^4$. For the projective-plane cycle and the
cycle from $H_4(N_{f_0})$ the matrix is
\[
 Q_W=\begin{pmatrix}0&1\\1&0\end{pmatrix}. \tag{P3F.22.3}
\]
The characteristic class $z$ extends $x$ and is positive on the
projective line. With the source's relative classes $u,v$ and natural
map $j$, its restrictions are $z^2=(1,0)$ and $p_1=(3,p)$, equivalently
$z^2=jv$ and $p_1=p\,ju+3jv$. The cited calculation therefore gives
\[
 \sigma(W)=0,\qquad p_1^2[W,\partial W]=6p,\qquad
 z^2p_1[W,\partial W]=p,\qquad z^4[W,\partial W]=0.
 \tag{P3F.22.4}
\]

\item The spin-c Eells--Kuiper formula
\cite[Section~3.1(v), pp.~5--6]{Jiang2014} applies: the boundary is a
homotopy sphere, so its rational degree-four relative-to-absolute map
is an isomorphism. It yields
\[
 \mu(-\lambda(f_0))=
 \frac{p_1^2-4\sigma(W)}{896}
 -\frac{z^2p_1}{192}+\frac{z^4}{384}
 =\frac{6p}{896}-\frac p{192}=\frac p{672}.
\]
Reverse the displayed boundary orientation and substitute $p=-24A(f)$:
\[
 \mu(\lambda(f))=\frac{A(f)}{28}\pmod{\mathbb Z}. \tag{P3F.22.1}
\]
The additive Eells--Kuiper classification identifies $\Theta_7$ with
$\frac1{28}\mathbb Z/\mathbb Z$, and the oriented Milnor sphere
$S_{2,-1}(\mathbb H)$ has invariant $+1/28$
\cite[Sections~2--3, pp.~67--69]{Davis1982}.
Jin's additivity of $A$ now proves additivity of $\lambda$. The pairs
$(k,\ell)(a)=(1,4)$ and $(k,\ell)(b)=(8,8)$ give
$\lambda(a)=0$, $\mu(\lambda(b))=1/28$ and
$\ker\lambda=\mathbb Za\oplus28\mathbb Zb$.
\end{enumerate}
\end{proof}
